% Cronograma do TCC. Cada linha é uma etapa; cada coluna, um mês (set/2025 a dez/2026).
% \cF = etapa realizada no mês; \cP = etapa prevista. Para ajustar, mova os \cF/\cP entre as colunas.
\definecolor{cronoFeito}{rgb}{0.255,0.345,0.176}
\definecolor{cronoPrevisto}{gray}{0.75}
\newcommand{\cF}{\textcolor{cronoFeito}{\rule{\linewidth}{1.6ex}}}
\newcommand{\cP}{\textcolor{cronoPrevisto}{\rule{\linewidth}{1.6ex}}}

\begin{table}[H]
    \centering
    \caption[Cronograma de execução do trabalho]{Cronograma de execução do trabalho, de setembro de 2025 a dezembro de 2026. Cada linha é uma etapa e cada coluna é um mês, identificado pela inicial (S = setembro, O = outubro, N = novembro, D = dezembro, J = janeiro e assim por diante). Células verdes indicam os meses em que a etapa foi realizada; células cinza, os meses em que ela está prevista. O primeiro bloco corresponde ao TCC1, dedicado ao planejamento e à elaboração do projeto, e o segundo ao TCC2, dedicado à execução técnica e à escrita da monografia.}
    \label{tab:cronograma}
    \scriptsize
    \setlength{\tabcolsep}{1.6pt}
    \renewcommand{\arraystretch}{1.15}
    \begin{tabular}{>{\raggedright\arraybackslash}p{6.6cm} *{16}{>{\centering\arraybackslash}p{0.42cm}}}
        \toprule
        & \multicolumn{4}{c}{\textbf{2025}} & \multicolumn{12}{c}{\textbf{2026}} \\
        \cmidrule(lr){2-5} \cmidrule(lr){6-17}
        \textbf{Etapa} & S & O & N & D & J & F & M & A & M & J & J & A & S & O & N & D \\
        \multicolumn{17}{l}{\textbf{TCC1: planejamento e projeto}} \\
        Background e fundamentação & \cF &  &  &  &  &  &  &  &  &  &  &  &  &  &  &  \\
        Revisão bibliográfica: gestão fiscal & \cF & \cF &  &  &  &  &  &  &  &  &  &  &  &  &  &  \\
        Revisão bibliográfica: legal (LGPD) & \cF & \cF &  &  &  &  &  &  &  &  &  &  &  &  &  &  \\
        Revisão bibliográfica: técnica (XAI) & \cF & \cF &  &  &  &  &  &  &  &  &  &  &  &  &  &  \\
        Problematização e objetivos &  & \cF &  &  &  &  &  &  &  &  &  &  &  &  &  &  \\
        Elaboração do projeto &  & \cF & \cF &  &  &  &  &  &  &  &  &  &  &  &  &  \\
        Escrita: referencial teórico &  & \cF & \cF &  &  &  &  &  &  &  &  &  &  &  &  &  \\
        Escrita: introdução e justificativa &  & \cF & \cF &  &  &  &  &  &  &  &  &  &  &  &  &  \\
        Escrita: metodologia &  &  & \cF & \cF &  &  &  &  &  &  &  &  &  &  &  &  \\
        Análise &  &  & \cF & \cF &  &  &  &  &  &  &  &  &  &  &  &  \\
        Cronograma &  &  & \cF &  &  &  &  &  &  &  &  &  &  &  &  &  \\
        Defesa do TCC1 &  &  &  & \cF &  &  &  &  &  &  &  &  &  &  &  &  \\
        \midrule
        \multicolumn{17}{l}{\textbf{TCC2: execução e monografia}} \\
        Reprodução do pipeline de produção &  &  &  &  & \cF & \cF & \cF &  &  &  &  &  &  &  &  &  \\
        Geração e validação dos dados sintéticos &  &  &  &  &  &  & \cF & \cF & \cF &  &  &  &  &  &  &  \\
        Aplicação do SHAP (global e local) &  &  &  &  &  &  &  & \cF & \cF & \cF &  &  &  &  &  &  \\
        Aplicação do LIME e análise de estabilidade &  &  &  &  &  &  &  &  & \cF & \cF &  &  &  &  &  &  \\
        Comparação entre SHAP e LIME &  &  &  &  &  &  &  &  &  & \cF & \cF &  &  &  &  &  \\
        Estrutura visual de apoio ao auditor &  &  &  &  &  &  &  &  & \cF & \cF &  &  &  &  &  &  \\
        Escrita da monografia &  &  &  &  &  &  &  &  &  & \cF & \cF & \cF & \cF &  &  &  \\
        Ajustes solicitados pelo orientador &  &  &  &  &  &  &  &  &  &  &  &  &  & \cF & \cF &  \\
        Experimento de desbalanceamento &  &  &  &  &  &  &  &  &  &  &  &  &  &  & \cF &  \\
        Revisão final e entrega &  &  &  &  &  &  &  &  &  &  &  &  &  &  & \cP &  \\
        Defesa do TCC2 &  &  &  &  &  &  &  &  &  &  &  &  &  &  &  & \cP \\
        \bottomrule
    \end{tabular}
\end{table}
